\chapter{Les stirges}
\sessioninfo{Session 12}{Mardi 24 novembre 2026}

Le deuxième palier de Firefinger s'enfonçait dans l'obscurité. Sous les vêtements de Destos, la lumière de son symbole sacré dessinait un halo discret sur le sol et révélait deux passages sinueux, l'un tourné vers l'est, l'autre vers le sud. Un vieux poteau brisé prit un instant l'apparence prometteuse d'un coffre avant que Morrÿs ne reconnaisse de simples morceaux de bois. Dans cette caverne, même les ombres semblaient décidées à se jouer d'eux.

Le nain posa une main contre la roche et étendit de nouveau ses sens à travers la pierre. Il ne perçut d'abord que les déplacements de ses compagnons. Le groupe choisit alors le couloir du sud et s'y engagea sans se séparer. À peine plusieurs d'entre eux eurent-ils pénétré dans la salle suivante que des vibrations s'éveillèrent au-dessus de Morrÿs.

Il leva la tête et avertit les autres. Neuf créatures se détachèrent du plafond dans un battement d'ailes désordonné. Elles avaient la taille d'un poing, un corps velu porté par des ailes membraneuses et une longue trompe effilée. Elles évoquaient des moustiques auxquels la jungle aurait accordé beaucoup trop d'ambition.

\illustration[0.8\textwidth]{2026-11-24/strige_portrait.png}
{\textit{Les stirges fondirent du plafond, leur longue trompe déjà dirigée vers leurs proies.}}
{stirge_portrait}

Azaka les accueillit d'un regard contrarié.

\begin{quote}
--- Encore des merdouilles.
\end{quote}

\scenebreak

Destos chercha l'origine des cris et ne comprit la menace qu'après que Morrÿs lui eut désigné le plafond. Une flamme sacrée jaillit vers la nuée, sans atteindre sa cible. Azaka arma aussitôt son arc et abattit l'une des créatures. Morrÿs, trop proche pour employer confortablement sa magie, tira sa dague et entailla une stirge qui continua pourtant de voleter autour de lui.

Kael posa sa marque sur une autre créature et la tua d'un coup d'épée. Les survivantes se laissèrent alors tomber sur le groupe. L'une d'elles manqua Morrÿs, mais une seconde planta sa trompe dans son épaule et demeura suspendue à lui. Une autre s'accrocha à Kael, tandis qu'une troisième atteignait Azaka.

La plaie de la guide ne saigna pas. Une fois encore, ses compagnons virent sa chair se refermer presque aussitôt autour de la blessure. Ils n'eurent cependant pas le temps de s'interroger : d'autres stirges voletaient déjà entre eux. Kael en abattit une qui tentait de gagner Destos, puis Obélik s'approcha d'Azaka, lui demanda de ne pas bouger et trancha proprement la créature fixée sur elle.

Azaka dégaina son cimeterre et tua la stirge accrochée à Kael. Morrÿs tenta de se débarrasser de celle qui lui pompait le sang, sans parvenir à l'atteindre. Kael voulut l'aider, mais retint son geste par crainte de frapper le nain. Pendant ce temps, la créature continuait de se gorger.

Un liquide noirâtre commença à suinter de la blessure de Morrÿs. Ses jambes cédèrent et il s'effondra, inconscient. Une autre stirge s'enfonça dans les chairs de Destos, tandis qu'une nouvelle attaque atteignait Kael. L'orc tomba à son tour.

\illustration[0.95\textwidth]{2026-11-24/attaque-stirges.png}
{\textit{Dans l'étroite caverne, les stirges se jetèrent sur les aventuriers pour leur voler leur sang.}}
{attaque_stirges_firefinger}

Obélik se plaça entre ses deux compagnons. D'un coup de hache, il fit éclater la créature accrochée à Morrÿs dans un mélange de sang et de liquide sombre. Destos arracha celle qui le piquait et la projeta au loin. Azaka, après avoir laissé échapper son arme, écrasa une autre stirge entre ses paumes comme un vulgaire insecte.

La créature fixée à Morrÿs abandonna le corps immobile pour chercher une proie plus vigoureuse. Elle se jeta sur Obélik et planta sa trompe dans son flanc. Le goliath tenta de l'atteindre, mais sa hache passa à côté. Destos se pencha alors sur Morrÿs et essaya d'arrêter son déclin, sans succès.

Azaka acheva la dernière stirge à mains nues. Elle se précipita ensuite vers les deux corps étendus et parvint à stabiliser Morrÿs, puis Kael. Le silence revint dans la salle, seulement troublé par les respirations lourdes des survivants. Une dizaine de minutes plus tard, les deux blessés rouvrirent enfin les yeux.

\begin{narrateur}
Les stirges avaient au moins une forme d'éthique : elles ne suçaient pas les morts. Elles tenaient simplement à vérifier de très près qu'ils l'étaient vraiment.
\end{narrateur}

\scenebreak

Azaka observa leurs visages épuisés et imposa une pause. Elle se plaça à l'entrée de la salle pour monter la garde pendant que les autres soignaient leurs blessures et reprenaient des forces. L'heure qui suivit permit au groupe de retrouver un peu de vigueur. Tous restaient néanmoins marqués par la violence d'un affrontement livré contre des créatures à peine plus grandes qu'un oiseau.

En s'installant, ils découvrirent dans un coin le cadavre exsangue d'un explorateur. Les stirges l'avaient entièrement vidé de son sang. Morrÿs fouilla ses affaires et trouva une bourse contenant quinze pièces d'or, un anneau en or et deux petites figurines taillées dans l'onyx. Chacune représentait un être à la tête ronde, percée de deux petits yeux.

Le nain montra ses trouvailles à ses compagnons. Azaka reconnut dans les statuettes des représentations de chwingas, de petits esprits de la forêt autour desquels circulaient de nombreuses légendes. Elle n'en révéla pas davantage. Morrÿs rangea soigneusement les figurines dans son sac avec le reste du butin.

La guide estimait qu'ils n'étaient plus très loin du sommet. Revenir plus tard signifierait devoir franchir une nouvelle fois tous les obstacles laissés derrière eux. Le groupe choisit donc de continuer malgré la fatigue. Destos conserva ses forces pour les dangers à venir et Azaka promit de ne laisser tomber personne, à condition que les plus mal en point sachent rester derrière elle.

\scenebreak

Un courant d'air provenait du sud. Le tunnel débouchait sur l'extérieur, où une corniche large d'à peine trente centimètres contournait la paroi de Firefinger. Cinquante mètres plus bas, la jungle s'étendait dans l'obscurité. Les bourrasques frappaient la tour de manière irrégulière, tantôt simples souffles, tantôt poussées assez violentes pour menacer l'équilibre de quiconque s'aventurerait dehors.

\illustration[0.95\textwidth]{2026-11-24/corniche.png}
{\textit{Balayée par les bourrasques, l'étroite corniche contournait Firefinger à cinquante mètres au-dessus de la jungle.}}
{corniche_firefinger}

Une échelle descendait le long de la roche, mais ses barreaux disparaissaient après une dizaine de mètres. Toute retraite par ce chemin était impossible. Morrÿs parcourut la salle avec sa baguette des secrets, à la recherche d'une trappe ou d'un passage dissimulé. Il ne trouva rien. Kael explora également le couloir de l'est, qui ne menait qu'à un cul-de-sac.

Il ne restait que la corniche. Vers l'est, elle se rétrécissait rapidement jusqu'à presque disparaître. Vers l'ouest, elle semblait conserver une largeur à peu près praticable. Le groupe choisit donc cette direction.

Obélik attacha une corde autour de Kael, puis assura l'autre extrémité autour de lui. Le rôdeur s'engagea lentement sur la corniche, le corps plaqué contre la paroi. Sous ses pieds, le vide descendait jusqu'à la jungle ; autour de lui, le vent cherchait sans cesse à le repousser. La roche était rugueuse, mais n'offrait aucune racine ni aucune prise à laquelle fixer une sécurité supplémentaire.

Kael progressa pas à pas jusqu'à apercevoir une structure de bois dressée contre la paroi. Il reconnut bientôt une nouvelle échelle, qui s'élevait sur environ vingt-cinq mètres vers le palier suivant. Arrivé à son pied, il fixa la corde d'Obélik à l'un de ses barreaux. Le passage disposait désormais d'une ligne tendue entre les deux extrémités de la corniche.

Restait à prévenir les autres. Kael imprima quelques secousses à la corde. Obélik sentit le signal, comprit que son compagnon cherchait à lui transmettre un message et répondit en tirant plus vigoureusement, sans saisir ce qu'on attendait de lui. Faute de langage plus précis, Kael s'assura avec sa propre corde et reprit prudemment le chemin du retour.

Une fois réuni au groupe, il leur exposa ce qu'il avait découvert. Ils fixèrent alors l'autre extrémité de la corde à l'échelle brisée. Désormais tendue entre les deux échelles, elle formerait une main courante le long de la corniche et leur offrirait également une voie de retraite. Chacun pourrait ainsi traverser à son tour sans qu'Obélik ait à demeurer en arrière pour assurer les autres.

Le dispositif était prêt. Devant eux, la corniche attendait le groupe, puis une nouvelle ascension de vingt-cinq mètres. Ils avaient survécu aux stirges ; il leur restait maintenant à vaincre le vent et le vide.
